\lecture{Sistemi di Trasporto Interni Automatici}{26~-~01~-~2026}
%
%
\section{Fabbrica Automatica e Sistemi Flessibili}
L'obiettivo della fabbrica automatica è l'integrazione del flusso dei materiali e delle informazioni per ottimizzare il processo produttivo tramite l'automazione industriale. La massima espressione di questo concetto è il Flexible Manufacturing System (FMS). Un sistema FMS è composto da:
\begin{itemize}
    \item Celle di fabbricazione: Ogni cella include un robot che serve diverse macchine utensili CNC o sistemi di saldatura;
    \item Sistema centrale di trasferimento: Collega le celle tramite nastri, catene o veicoli automatici (AGV);
    \item Calcolatore centrale: Coordina tutte le attività del sistema.
\end{itemize}

%
\section{Automated Guided Vehicles (AGV)}
Gli AGV sono veicoli per il trasporto interno senza guidatore a bordo. Un sistema AGV si compone di quattro elementi fondamentali:
\begin{enumerate}
    \item I carrelli: Alimentati a batterie, eseguono il trasporto fisico;
    \item Sistema di guida: Guida i carrelli lungo i percorsi stabiliti;
    \item Impianto per comunicazioni: Trasmette comandi e informazioni tra sistema di gestione e carrelli;
    \item Calcolatore di gestione: Controlla l'intero impianto.
\end{enumerate}

%
\subsection{Tecnologie di Guida}
Esistono diverse tecnologie per guidare questi veicoli:
\begin{itemize}
    \item \textbf{Guida Induttiva (Fissa)}: Un cavo interrato nel pavimento genera un campo magnetico rilevato da bobine sul carrello. Può essere monofrequenza (un solo cavo per tutti i percorsi) o multifrequenza (fili diversi con frequenze diverse per ogni percorso).
    \item \textbf{Guida Ottica (Fissa)}: Utilizza nastri catarifrangenti o strisce colorate a terra. I sensori possono essere a luce riflessa o telecamere a bordo.
    \item \textbf{Guida Laser (LGV)}: Navigazione odometrica con scansione laser. Una testa rotante sul veicolo invia raggi laser verso riflettori installati nell'ambiente. Il computer calcola la posizione triangolando almeno 3 o 4 punti di riflessione.
    \item \textbf{Guida Libera Visuale}: Telecamere stereo a 360° costruiscono una mappa 3D dinamica dell'ambiente (tecnologia Evidence Grid), eliminando la necessità di infrastrutture fisse.
    \item \textbf{Navigazione Stimata (Dead Reckoning)}: Software che calcola la posizione basandosi sui sensori di spostamento degli assi (encoder). Adatto solo per tragitti brevi poiché perde precisione sulla lunga distanza.
\end{itemize}

%
\subsection{Tipologie di carrelli}
Sono presenti diverse tipologie di carrelli:
\begin{itemize}
    \item \textbf{Carrelli trainanti (trattori)}: Trainano una serie di rimorchi, ideali per traffico elevato;
    \item \textbf{Carrelli trasportatori}: Dotati di piattaforme, rulli motorizzati o forche per il carico automatico;
    \item \textbf{Muletti a guida automatica}: Simili ai carrelli elevatori tradizionali ma senza operatore (es. LGV a forche);
    \item \textbf{Assembly Platforms}: L'AGV funge da postazione di lavoro mobile integrata nella linea di assemblaggio.
\end{itemize}

%
\subsection{Sicurezza e Progettazione}
La sicurezza è garantita da normative (UNI EN 1525) e dispositivi come laser scanner a 360 gradi, paraurti sensibili, pulsanti di emergenza e sensori che distinguono ostacoli fissi da quelli dinamici. La progettazione richiede l'analisi dei flussi, il calcolo del numero di veicoli (metodo deterministico), la definizione del layout e delle aree di ricarica.

%
\paragraph{Vantaggi} Flessibilità, riduzione scorte e danni, sicurezza, integrazione con altri sistemi.
%
\paragraph{Svantaggi} Alti costi d'investimento, sensibilità ambientale (pavimenti, clima) e inoltre necessità di manutenzione specializzata.

%
\section{Robot Industriali}
Il robot è un manipolatore multifunzionale riprogrammabile progettato per muovere materiali o attrezzi attraverso movimenti programmati. Componenti principali:
\begin{itemize}
    \item Manipolatore: La struttura meccanica composta da bracci (link) e giunti;
    \item Attuatori: Motori (elettrici, idraulici, pneumatici) che generano il movimento;
    \item End Effector: L'organo terminale (pinza, saldatrice) collegato alla flangia finale;
    \item Controllo: Il cervello elettronico che gestisce traiettorie, I/O e interazione.
\end{itemize}

%
\paragraph{Gradi di libertà} Rappresentano i movimenti consentiti. Per posizionare e orientare un oggetto in qualsiasi punto dello spazio 3D sono necessari 6 gradi di libertà (3 per la posizione, 3 per l'orientamento).

%
\paragraph{Sensori} I sensori si suddividono in:
\begin{itemize}
    \item \emph{Sensori per contatto}: Tattili, di forza/coppia.
    \item \emph{Sensori senza contatto}: Visivi, di prossimità, laser.
\end{itemize}

%
\paragraph{Applicazioni} Saldatura (punti/arco), verniciatura, assemblaggio, carico/scarico macchine, pallettizzazione.

%
\section{Industria 4.0 e Lean}
L'Industria 4.0 rappresenta la quarta rivoluzione industriale, basata sui Sistemi Cyber-Fisici dove oggetti reali e processi virtuali sono interconnessi. Il concetto chiave è la connessione di persone, macchine e sistemi per creare una Fabbrica Intelligente, capace di adattarsi autonomamente. Le tecnologie abilitanti includono:
\begin{itemize}
    \item Internet of Things (IoT);
    \item Big Data e Analytics;
    \item Cloud Manufacturing;
    \item Additive Manufacturing.
\end{itemize}
Il \emph{Lean Management} è propedeutico a Industria 4.0: l'approccio snello e il coinvolgimento dei dipendenti creano le condizioni per il successo dell'automazione.

%
\section{Robot Collaborativi}
A differenza dei robot tradizionali (che lavorano in gabbie), i CoBot sono progettati per interagire fisicamente con l'uomo nello stesso spazio di lavoro, senza barriere.
%
\paragraph{Caratteristiche distintive}
\begin{itemize}
    \item Più lenti e meno potenti dei robot industriali per garantire la sicurezza;
    \item Facili da programmare (apprendimento guidato dall'operatore) e riposizionare;
    \item Non sostituiscono l'uomo, ma lo sollevano da compiti ripetitivi, faticosi o a basso valore aggiunto.
\end{itemize}

%
\paragraph{Integrazione Mobile (AMR)} Un CoBot montato su un AGV diventa un Autonomous Mobile Robot. Questo sistema può navigare autonomamente tra gli scaffali, eseguire picking o sanificazione degli ambienti, lavorando giorno e notte.

%
\section{Logistica 4.0}
Applicazione dei paradigmi di Industria 4.0 alla Supply Chain. Si basa su tre parole chiave: automazione, interconnessione e integrazione. Le tecnologie principali applicate alla logistica sono:
\begin{itemize}
    \item Sensori IoT: Per tracking in tempo reale e monitoraggio condizioni (es. temperatura cibi);
    \item Cloud: Per la condivisione sicura dei dati tra gli attori della filiera;
    \item Analytics: Per estrarre valore dai dati e ottimizzare i flussi;
    \item Realtà Aumentata: Supporto al picking per ridurre errori;
    \item Additive Manufacturing: Utile per la gestione dei ricambi on-demand (postponement logistico).
\end{itemize}